\section{Minor geometry and rational certificates}\label{mi:appendix}

\subsection{Orbit, support and the tangent pencil}

\begin{proof}[Proof of \cref{mi:orbit}]\label{mi:orbit-proof}
Write $M=H+tJ$, with $H$ symmetric. Then
$\det M=\det H+t^2$. Thus
$C_I=\{H+tJ:H\succeq0,t\in\R\}$ is closed, convex and free, with
interior $H\succ0$.
To prove maximality, take $N\notin C_I$ and diagonalize its symmetric
part. It has a negative eigenvalue. Choose a positive definite symmetric
matrix $H$ small enough in that eigendirection and large enough in the
orthogonal direction that $\sym N+H$ is indefinite. Set
$Z=H-\tfrac12(N-N^T)\in\intt C_I$.
Then $(N+Z)/2$ is symmetric and has negative determinant.
If a convex set contained both $C_I$ and $N$, this midpoint would be an
interior point: a ball around $Z$ remains in $C_I$, and its midpoints
with $N$ form a ball. Its interior also contains $I$, with positive
determinant. The segment between these interior points contains a
singular matrix in the interior, contradicting freeness. This proves
maximality without a classification theorem.

For $N=AMB^T$, positive semidefiniteness of
$\sym(F^TA^{-1}NB^{-T})$ is equivalent, by congruence, to that of
$\sym(BF^TA^{-1}N)$. This gives $F'^T=BF^TA^{-1}$.
Taking $A=I$, $B=F^T$ shows that all $\det F>0$ occur in the orbit of
$C_I$; maximality transfers under invertible transformations.
The substitution $u=Fv$ in $v^TMFv\ge0$ gives the transposition rule.
Moreover, with the matrix trace inner product,
$C_F^*=F\mathbb S_+^2$. Equality of two cones implies that
$G=F^{-1}F'$ maps $\mathbb S_+^2$ onto itself by left multiplication.
Since symmetric matrices are spanned by positive semidefinite matrices,
$G$ preserves the symmetric matrix space. Testing $I,E_{11},E_{22}$,
and $E_{12}+E_{21}$ gives $G=cI$; preservation of the positive cone gives
$c>0$. Hence the projective family has dimension three.

Let $A\bar MB^T=U'H'$ be a transformed polar decomposition.
Pulling $C_{U'}$ back gives
$F^T=B^{-1}U'^TA$ and
$F^T\bar M=B^{-1}H'B^{-T}\succ0$ symmetric.
Conversely, if $F^T\bar M=H\succ0$, apply the coordinate map $M\mapsto F^TM$;
the transformed polar factor is $I$, so its pullback is $C_F$.
This family is parametrized by $H\succ0$, modulo a positive scalar,
and has dimension two. The parameter transformation sends $H$ to
$BHB^T$, proving equivariance; for transposition it sends $H$ to
$F^{-T}HF^{-1}$.
Finally $R^T\bar M\succ0$ symmetric with $R\in SO(2)$ forces $R=U$ by
polar uniqueness. The negative polar sign gives the opposite cone and is
not a second point-rule member.
\end{proof}

\begin{proof}[Proof of \cref{mi:support}]\label{mi:support-proof}
Positive objective weights make finite objective sublevels compact, so
the minimum is attained. Choose a minimizer with smallest support.
If its supported rays are dependent, choose a nonzero supported vector
$h$ with $Ph=0$. Both small signs of $h$ preserve nonnegativity and the
same constraint point; optimality forces $\rc^Th=0$. Move until a
coordinate first vanishes, contradicting smallest support. Thus an
independent-supported minimizer exists.

Fix any independent-supported minimizer with support $J$ of size $k$.
Put $W=\operatorname{span}\{P_j:j\in J\}$ and
$V=P_J\{h:\rc_J^Th=0\}$, so $\dim V=k-1$.
The determinant is nonnegative on $T_{\zK}$: a negative value there
would produce a lower-cost zero on the segment from the apex.
Small directions in $V$ remain in the supported objective-level face,
so $B(M_*,V)=0$ and the quadratic form is nonnegative on $V$.
If the determinant gradient vanishes on $W$, the same argument for
directions with decreasing objective, together with evenness of the
quadratic term, makes the form nonnegative on all of $W$, so $k\le2$.
Otherwise $V$ is a nonnegative subspace of dimension $k-1$ of a form
with positive index two; hence $k\le3$.

If $k=3$, the preceding exceptional case is impossible, and $M_*\ne0$.
The restriction of the form to
$T=\{D_{\mathrm{edge}}:B(M_*,D_{\mathrm{edge}})=0\}$ has radical $\R M_*$ and signature $(1,1)$ after
quotienting by that radical. Indeed, choose $Z$ with $B(M_*,Z)=1$;
$\operatorname{span}\{M_*,Z\}$ has signature $(1,1)$ and its orthogonal
complement has signature $(1,1)$.
The two-dimensional nonnegative subspace $V\subseteq T$ must contain
$M_*$, since otherwise its image in the quotient would be a
two-dimensional nonnegative subspace of signature $(1,1)$.
The supported objective-level affine plane $M_*+V$ therefore contains
zero. As it lies in $\bar M+W$, this implies $\bar M\in W$.
For $k=2$, local nonnegativity along its edge gives
$B(M_*,D_{\mathrm{edge}})=0$ and $\det D_{\mathrm{edge}}\ge0$, as asserted.
\end{proof}

\begin{proof}[Proof of \cref{mi:pencil}]\label{mi:pencil-proof}
The two-dimensional identities
\[
 \adj(ab^T)=(Jb)(Ja)^T,\qquad \adj(D_{\mathrm{edge}})=J^TD_{\mathrm{edge}}^TJ
\]
give $G_1a=0$ and
$(Jb)^TG_0a=b^T\adj(D_{\mathrm{edge}})a=\tr(\adj(M_*)D_{\mathrm{edge}})=0$.
Thus $G_0a=c_0b$; invertibility of $G_0$ gives $c_0\ne0$.
Linearity of the adjugate and tangency give
\[
 \det(D_{\mathrm{edge}}+tM_*)=\det D_{\mathrm{edge}},
 \qquad J\adj(D_{\mathrm{edge}}+tM_*)=G_0+tG_1,
 \qquad\det(G_0+tG_1)=\det D_{\mathrm{edge}}.
\]
At the contact point, $\sym((F^Ta)b^T)\succeq0$ forces $F^Ta=cb$ for
$c>0$: if the two vectors are not parallel the symmetric matrix is
indefinite, and invertibility rules out $c=0$.
Write $Y=\sym(F^TD_{\mathrm{edge}})$ and $u=Jb$.
The segment condition is $cbb^T+sY\succeq0$ for both signs of $s$.
It first gives $u^TYu=0$; positive semidefiniteness then gives
$(cbb^T+sY)u=0$, hence $Yu=0$.
Therefore $Y=\mu bb^T$, and with $t=-\mu/c$,
$F^T(D_{\mathrm{edge}}+tM_*)$ is skew-symmetric. It equals $\eta J$, so
\[
 F^T=\frac{\eta}{\det D_{\mathrm{edge}}}J\adj(D_{\mathrm{edge}}+tM_*)
     =\theta(G_0+tG_1).
\]
The contact identity $F^Ta=cb$ fixes $\theta c_0>0$.
\end{proof}

\subsection{Complete data and proofs for A, B and S1}

All entries below use the order $(a,b,c,d)$. Put $V_0=\bar M$,
$P_j=V_j-V_0$ and $\rc_j=1$.
\begin{table}[htbp]
\centering\small
\begin{tabular}{c c c c}
\toprule
 & A & B & S1\\
\midrule
$V_0$ & $(0,5/2,-1/2,-1/2)$ & $(-3,-5/2,1/2,-2)$ & $(-4,-1,-1/2,-2)$\\
$V_1$ & $(-5,3,4,-4)$ & $(1,1,-7,-4)$ & $(3,3,-3,-3)$\\
$V_2$ & $(-8,12,-8,8)$ & $(-1/2,-1/2,-11/2,-7)$ & $(95/24,4,-4,-4)$\\
$V_3$ & $(-3/2,-4,2,-7/2)$ & $(1/2,1,-5/2,0)$ & $(17/3,6,-3/2,-3/2)$\\
$V_4$ & $(-7/2,-1/2,3/2,-4)$ & $(-7/2,4,-7/2,2)$ & $(41/6,3,-6,-2)$\\
$M_*$ & $(-6,6,0,0)$ & $(0,0,-6,-6)$ & $(3,3,-3,-3)$\\
$\lambda_*$ & $(2/3,1/3,0,0)$ & $(1/3,2/3,0,0)$ & $(1,0,0,0)$\\
$\det\bar M$ & $5/4$ & $29/4$ & $15/2$\\
$\det[P_1\ \cdots\ P_4]$ & $-2775/4$ & $1005/4$ & $533/16$\\
\bottomrule
\end{tabular}
\caption{Rational full-dimensional minor corners.}\label{mi:example-data}
\end{table}

\paragraph{Exact minimization on a simplex.}
For $V(\beta)=\sum_{i=0}^4\beta_iV_i$, let
$H_{ij}=B(V_i,V_j)$. Then $\det V(\beta)=\beta^TH\beta$ on
$\beta\ge0$, $\mathbf1^T\beta=1$.
On every nonempty face $I$, an interior stationary point solves
\begin{equation}\label{mi:face-system}
 \begin{pmatrix}2H_{II}&-\mathbf1\\\mathbf1^T&0\end{pmatrix}
 \begin{pmatrix}\beta_I\\\mu\end{pmatrix}
 =\begin{pmatrix}0\\1\end{pmatrix}.
\end{equation}
Only solutions with all barycentric coordinates positive need be retained.
A singular system cannot hide a new minimum: when a stationary solution
exists, a null vector $(h,\eta)$ satisfies
$2H_{II}h=\eta\mathbf1$, $\mathbf1^Th=0$.
Multiplying by the stationary barycentric vector gives $\eta=0$.
Thus the objective is constant along $\beta_I+th$, and its minimum
also occurs on lower-dimensional faces. A nonconstant segment of minima
would have two distinct boundary minimizers. This proves that face
enumeration establishes both the value and uniqueness.

For A, B and S1, the respective polar matrices are
\begingroup\small
\[
H_{\mathrm A}=\begin{pmatrix}
5/4&-3&15&-25/8&-9/8\\
-3&8&-16&67/4&63/4\\
15&-16&32&-20&-9\\
-25/8&67/4&-20&53/4&101/8\\
-9/8&63/4&-9&101/8&59/4
\end{pmatrix},\qquad
H_{\mathrm B}=\begin{pmatrix}
29/4&-4&17/4&-31/8&-39/8\\
-4&3&-3/2&15/4&95/4\\
17/4&-3/2&3/4&3/8&175/8\\
-31/8&15/4&3/8&5/2&29/4\\
-39/8&95/4&175/8&29/4&7
\end{pmatrix},
\]
\[
H_{\mathrm{S1}}=\begin{pmatrix}
15/2&9/4&73/24&-23/12&-61/12\\
9/4&0&1/16&1/2&1/4\\
73/24&1/16&1/6&67/96&3/8\\
-23/12&1/2&67/96&1/2&227/24\\
-61/12&1/4&3/8&227/24&13/3
\end{pmatrix}.
\]
\endgroup
All 31 systems in each example are nonsingular. The following table
lists every stationary value whose barycentric coordinates are positive;
``--'' means that no such stationary point occurs. A face is named by
its vertex indices. The five vertex values are the corresponding diagonal
entries of $H$.
\begin{center}\small
\begin{tabular}{c r r r}
\toprule
Face & A & B & S1\\\midrule
01&$4/61$&$23/73$&--\\
03&$435/1328$&$199/1120$&$11/1704$\\
04&$1099/1168$&$1727/1536$&$959/3168$\\
12&$0$&$0$&--\\
13&$57/4$&$105/32$&--\\
14&$2081/140$&$8689/600$&--\\
23&$96/341$&$111/160$&$3721/6720$\\
24&$1564/259$&$30289/2304$&--\\
34&$2307/176$&$561/80$&$50281/8112$\\
012&$36/37$&$27/70$&--\\
013&$6497/16997$&$1048/2361$&$2511/6704$\\
014&--&$204432/73729$&$162/265$\\
023&$88880/55709$&$7193/7848$&$33595/64082$\\
024&--&$573681/114824$&$1487/1829$\\
034&--&$2317/1929$&$575117/361346$\\
123&$864/619$&--&--\\
134&$30431/2165$&--&--\\
0123&$142884/67651$&--&--\\
0234&--&--&$15689521/9829632$\\
\bottomrule
\end{tabular}
\end{center}
Every unlisted face of size at least two has no positive-barycentric
stationary point. Thus the unique zero is on face 12 with weights
$(2/3,1/3)$ for A and $(1/3,2/3)$ for B, and at vertex 1 for S1.
It follows that $\zK=1$ with the stated unique minimizers.

\begin{proof}[Remaining assertions of \cref{mi:examples}]\label{mi:example-proof}
For A, $D_{\mathrm{edge}}=V_1-V_2=(3,-9,12,-12)$ has determinant 72 and
$\nabla\det(M_*)^TD_{\mathrm{edge}}=0$. Its pencil is
\[
 G_0=\begin{pmatrix}-12&3\\12&-9\end{pmatrix},\qquad
 G_1=\begin{pmatrix}0&-6\\0&6\end{pmatrix}.
\]
The positive sign of the pencil scale is required at $M_*$.
The intervals on its parameter imposed by the five vertices are
\begin{align*}
 K(V_0)&=[6-\sqrt{10},6+\sqrt{10}],&
 K(V_1)&=(-\infty,3],&K(V_2)&=[-3/2,\infty),\\
 K(V_3)&=[-49/6-2\sqrt{106}/3,-49/6+2\sqrt{106}/3],\\
 K(V_4)&=[-33/10-2\sqrt{118}/5,-33/10+2\sqrt{118}/5].
\end{align*}
These follow by the two diagonal and determinant inequalities for
$\sym((G_0+tG_1)V_i)$. Already $K(V_0)$ and $K(V_3)$ are disjoint.
For B, $D_{\mathrm{edge}}=(3/2,3/2,-3/2,3)$ has determinant $27/4$ and its sign-normalized
pencil is
\[
 G_0=\begin{pmatrix}-3/2&-3/2\\3&-3/2\end{pmatrix},\qquad
 G_1=\begin{pmatrix}-6&0\\-6&0\end{pmatrix}.
\]
The intervals are
\begin{align*}
 K(V_0)&=[(11-\sqrt{87})/2,(11+\sqrt{87})/2],&
 K(V_1)&=(-\infty,3/2],&K(V_2)&=[-3,\infty),\\
 K(V_3)&=[-9/4-\sqrt{30}/2,-9/4+\sqrt{30}/2],\\
 K(V_4)&=[-1/12-\sqrt{21}/15,-1/12+\sqrt{21}/15].
\end{align*}
Again $K(V_0)$ and $K(V_3)$ are disjoint.
The multipliers in $\rc=-\sigma P^T\nabla\det(M_*)+\nu$ are
\[
\begin{array}{c|c|c|c}
 &\sigma&\nu&\nabla\det(M_*)^T\bar M\\\hline
\mathrm A&1/6&(0,0,3/2,5/2)&6\\
\mathrm B&1/3&(0,0,1,15)&3\\
\mathrm{S1}&2/9&(0,1/36,2/9,1/9)&9/2
\end{array}
\]
For S1, the normal products on the four incident edges to
$V_0,V_2,V_3,V_4$ are respectively $9/2,1/8,1,1/2$, all positive.
For A, the unique corner contact satisfies
$\nabla\det(M_*)^T(\bar M-M_*)=6>0$.
The attainment criterion in \cref{fd:attainment}, applied to
$\{\det\le0\}$ and then \cref{mi:free-cones}, gives an unrestricted
maximal free set attaining one.

For the strict supremum gaps, set
$W_0=V_0$ and $W_j=(1-z)V_0+zV_j$.
Represent a symmetric certificate matrix by
$[u,v,h]=\begin{psmallmatrix}u&v\\v&h\end{psmallmatrix}$.
The following complete rational certificates satisfy
$\sum_{i=0}^4W_iY_i=0$ and $Y_i\succ0$.
Both assertions are checked by matrix multiplication and by
$u>0$, $uh-v^2>0$ in each row.
\begin{center}\small
\begin{tabular}{c r r r}
\toprule
\multicolumn{4}{c}{A, $z=9/20$}\\
$i$&$u$&$v$&$h$\\\midrule
0&$120111307/200000000$&$22844639/100000000$&$8690461/100000000$\\
1&$2000019/2000000$&$28763/50000$&$330939/1000000$\\
2&$26271/1000000$&$6631/200000$&$20937/500000$\\
3&$43823/1000000$&$-25577/1000000$&$14943/1000000$\\
4&$11/1000000$&$0$&$11/1000000$\\\midrule
\multicolumn{4}{c}{B, $z=461/500$}\\
0&$56409442877/1189000000000$&$-400381181/3625000000$&$93214817/362500000$\\
1&$58606291/82000000$&$-211351/250000$&$1$\\
2&$15799/20000$&$-382359/500000$&$370147/500000$\\
3&$1/1000000$&$0$&$1/500000$\\
4&$4523/1000000$&$14917/1000000$&$49221/1000000$\\\midrule
\multicolumn{4}{c}{S1, $z=779/1000$}\\
0&$68915327/21600000000$&$3999047/1800000000$&$11344801/7200000000$\\
1&$407/21600000$&$-1/100000$&$1/50000$\\
2&$11101/100000$&$-12781/100000$&$3679/25000$\\
3&$1$&$-76729/100000$&$471/800$\\
4&$3653/100000$&$-6377/100000$&$11139/100000$\\
\bottomrule
\end{tabular}
\end{center}
\Cref{mi:certificates} now gives all three upper bounds.

For robustness, the linear map
$\mathcal L:(Y_i)\mapsto\sum_iW_iY_i$ from symmetric blocks to
$\R^{2\times2}$ has rank four. Its adjoint kernel consists of $F$ with
$\sym(F^TW_i)=0$ for every $i$.
The apex condition gives $F^T=tJV_0^{-1}$. If $t\ne0$, the other
conditions force every $V_0^{-1}W_i$ to be scalar, contradicting the
independence of the four rays. Hence the kernel is zero.
In a neighborhood of A, the map retains full row rank, so projecting the
given positive definite tuple onto its continuously varying kernel keeps
all blocks positive definite. The orbit upper bound $9/20$ persists.
The corner bound is continuous here. Lower semicontinuity follows from
compact objective sublevels and continuity of the determinant. For upper
semicontinuity, the feasible point $M_*+\varepsilon P_1$ has negative
determinant for small $\varepsilon>0$, since
$\nabla\det(M_*)^TP_1=-6$. This sign persists under small perturbations,
so the segment from the perturbed apex reaches a zero at cost arbitrarily
close to one. The strict gap therefore persists.
\end{proof}

The three coarse bounds are the portable certificates needed for the
counterexamples. Narrower archived verification intervals are reported
with the computational evidence; their numerical precision is not needed
for the preceding proofs.

\subsection{Portable comparisons and bilinear brackets}\label{mi:comparison-certificates}

Rational feasible parameters, checked at the indicated levels, include
\[
\begin{array}{c|c|c}
\text{corner}&z&F^T\\\hline
\mathrm A&449533/1000000&
\begin{pmatrix}-10188147307/2513&-24509003591/8528\\
34718214589/9897&-20333469614/8855\end{pmatrix}\\[6pt]
\mathrm B&921947/1000000&
\begin{pmatrix}-15696987146/8949&-650114959/308\\
2464903799/1862&-19805038891/9695\end{pmatrix}\\[6pt]
\mathrm{S1}&389053/500000&
\begin{pmatrix}2721203962/1605&-11697519049/801\\
47980238700/6149&-39926061217/9010\end{pmatrix}
\end{array}
\]
In each case $\det F>0$, $\sym(F^TV_0)\succ0$, and
$\sym(F^TW_i)\succeq0$ for $i=1,\ldots,4$.
Thus B's ratio exceeds $0.921947$, while its coarse upper bound is $0.922$.

For completeness, the three bilinear slices have the following data in
the same minor coordinate order. Their objective weights are all one.
\[
\begin{array}{c|c|c|c}
&\mathrm T&\mathrm T'&\mathrm V\\\hline
V_0&(3/2,-9/2,0,1)&(1,-1/2,1/2,1)&(2,-2,3,1)\\
V_1&(18,-1,-6,1)&(-1/2,1/2,-4,1)&(0,0,0,1)\\
V_2&(-18,-5,6,1)&(24,-10,3,1)&(1/4,6,-2,1)\\
V_3&(5/2,0,5/2,1)&(15/2,9/2,1/2,1)&(1/2,1,-5/2,1)\\
\lambda_*&(1/2,1/2,0)&(6/7,1/7,0)&(1,0,0)
\end{array}
\]
They are the two tangent-edge and support-one bilinear examples,
respectively. Their corner bounds are one by the corresponding bilinear
proofs. The following lower parameters prove the lower ends of the
brackets in \cref{mi:embedding} directly:
\begin{align*}
F_{\mathrm T}^T&=
\begin{pmatrix}3557730690827343/67108864&2666317851727117/16777216\\
-22689768773827/93440775&1900883772391943/99891300\end{pmatrix},
&z_{\mathrm T}^{\mathrm{lo}}&=9753853/10^7,\\
F_{\mathrm T'}^T&=
\begin{pmatrix}1151288814435793/80718478&-1212462267789092/97079703\\
-56281863190224/49263763&709421204462399/87313318\end{pmatrix},
&z_{\mathrm T'}^{\mathrm{lo}}&=8347663/10^7,\\
F_{\mathrm V}^T&=
\begin{pmatrix}722028544872305/33554432&-251095200246904/33842369\\
121179366244954/66838953&2714408638762620/85585871\end{pmatrix},
&z_{\mathrm V}^{\mathrm{lo}}&=1967693/2000000.
\end{align*}
The tests are the same rational determinant and diagonal tests as above.
At the respective upper levels
$z=1219233/1250000$, $2086943/2500000$, and $2459617/2500000$,
the following $Y_i=[u,v,h]$ give the complete dual certificates:
\begin{center}\scriptsize
\begin{tabular}{c r r r}
\toprule
\multicolumn{4}{c}{$\mathrm T$, $z=1219233/1250000$}\\
$i$&$u$&$v$&$h$\\\midrule
0&$284680208922551/1125000000000000$&$125759927996883/250000000000000$&$250000002793643/250000000000000$\\
1&$773793481/1800000000$&$15927467/100000000$&$5901223/100000000$\\
2&$34295509/100000000$&$-8578191/100000000$&$2145641/100000000$\\
3&$2439499/100000000$&$-509103/4000000$&$66403807/100000000$\\\midrule
\multicolumn{4}{c}{$\mathrm T'$, $z=2086943/2500000$}\\
0&$443012111728267/3750000000000000$&$141798875388169/1250000000000000$&$68080425260229/625000000000000$\\
1&$156624919/750000000$&$2856141/6250000$&$1$\\
2&$1033917/100000000$&$3925397/100000000$&$14903451/100000000$\\
3&$185279/50000000$&$-1727869/100000000$&$4028567/50000000$\\\midrule
\multicolumn{4}{c}{$\mathrm V$, $z=2459617/2500000$}\\
0&$283184926112931/2000000000000000$&$-243552272820433/2000000000000000$&$209466408183833/2000000000000000$\\
1&$1$&$41186779/400000000$&$1060221/100000000$\\
2&$7891043/50000000$&$-5452943/50000000$&$753629/10000000$\\
3&$287531/10000000$&$2685447/100000000$&$627031/25000000$\\
\bottomrule
\end{tabular}
\end{center}
Each block is positive definite and $\sum_{i=0}^3W_iY_i=0$.
These data prove both ends of all three brackets.
In particular,
$779/1000<8347663/10^7$ and
$8347772/10^7<921947/10^6<461/500<9753853/10^7$,
which prove every comparison in \cref{mi:embedding} without relying on
unprinted endpoint witnesses.

\subsection{Scaling and other minor projections}

\begin{proof}[Proof of \cref{mi:scaling}]\label{mi:scaling-proof}
The determinant is
\[
 \delta\bigl[(1-\lambda_2)(1-\lambda_3)+\lambda_1\lambda_4\bigr].
\]
It cannot vanish at cost below one, and it vanishes at $e_2,e_3$.
The polar factor is $I$. Its ray steps are
$(2\sqrt\delta,1,1,2/\sqrt\delta)$, by the determinant test on
$\sym(\bar M+tP_j)$. The point-rule parameter
$F^T=\diag(1,1/\delta)$ has $F^T\bar M=I$ and steps $(2,1,1,2)$.

For a rotation parameter
$F^T=\begin{psmallmatrix}c&s\\-s&c\end{psmallmatrix}$,
positive definiteness at the apex forces $c>0$ and
$|u|<2\sqrt\delta/(1-\delta)\le(8/3)\sqrt\delta$, where $u=s/c$.
The first ray's positive semidefiniteness requires
\[
 \bigl(t-u(1-\delta)\bigr)^2\le4(\delta-tu).
\]
For $u\ge0$ this implies $t<4\sqrt\delta$.
For $u<0$, expansion gives $t^2\le4\delta+4t|u|$, so
\[
 t\le2|u|+2\sqrt{u^2+\delta}
 <\frac{16+2\sqrt{73}}3\sqrt\delta.
\]
The cut bound is no greater than its first-ray step. Finally left
multiplication by $\diag(1,\delta)$ maps the stated unscaled corner to
this corner and maps its exact $C_I$ to the exact point-rule parameter;
the newly computed polar-factor choice is still $C_I$.
\end{proof}

\paragraph{Other minor projections.}\label{mi:other-minors}
For a principal minor, $M$ is symmetric and $\det M=ad-b^2$ has signature
$(1,2)$. The positive-determinant result in \cref{mi:principal} concerns
the full singular locus, including positive and negative rank-one matrices.
At a negative-determinant apex, reversing the form gives signature $(2,1)$;
the general homogeneous signature-$(n,1)$ orbit characterization applies
to that different component. It does not turn the polar point rule into
the full family. If the diagonal bounds are imposed, the forbidden
projection is instead the rank-one positive semidefinite cone. In that
setting, at an indefinite apex, the maximal outer-product-free sets
containing that apex in their interiors are halfspaces $\ip{A}{M}\ge0$ with
$A\preceq0$ and $\ip{A}{\bar M}>0$ \citep{BCM2020}.
Their intersection cuts are dominated on the corner cone by the valid
inequality $\ip{A}{M}\le0$. Indeed, writing
$-A=\mu_1g_1g_1^T+\mu_2g_2g_2^T$, $\mu_i\ge0$, shows this is a
nonnegative combination of at most two valid inequalities
$g_i^TMg_i\ge0$. A particular eigenvector separator need not select
these two vectors.

With exactly one diagonal entry, say
$a=x_i^2$, $b=x_ix_l$, $c=x_jx_i$, $d=x_jx_l$, the projection has closure
$\{\det M=0,a\ge0\}$.
For $a>0$, choose $x_i=\sqrt a$, $x_l=b/\sqrt a$, $x_j=c/\sqrt a$;
then $d=bc/a$ follows from the determinant equation.
For $a=0$, the determinant requires $bc=0$.
If $b\ne0$, approach with $a>0$, fixed $b,d$, and $c=ad/b$;
the case $c\ne0$ is symmetric. If $b=c=0$, choose $b,c$ tending to zero
with $bc=ad$, including either sign of $d$.
Thus every boundary point is a limit of projected outer products.
This smaller forbidden set can admit larger free sets and a larger corner
bound. The four-distinct-index orbit guarantees remain valid cuts, but
their ratios to that larger bound can only decrease.
